\documentclass[11pt,a4paper]{article}
\usepackage{amsmath,amsthm,amssymb,mathtools}
\usepackage[margin=2.5cm]{geometry}
\usepackage{hyperref}
\usepackage{enumitem}
\usepackage{booktabs}
\usepackage{xcolor}
\usepackage{tcolorbox}

\hypersetup{
  colorlinks=true,
  linkcolor=blue!70!black,
  citecolor=green!50!black,
  urlcolor=blue!60!black
}

% Theorem environments
\newtheorem{theorem}{Theorem}[section]
\newtheorem{lemma}[theorem]{Lemma}
\newtheorem{proposition}[theorem]{Proposition}
\newtheorem{corollary}[theorem]{Corollary}
\newtheorem{definition}[theorem]{Definition}
\newtheorem{example}[theorem]{Example}
\newtheorem{remark}[theorem]{Remark}
\newtheorem{conjecture}[theorem]{Conjecture}
\newtheorem{openproblem}{Open Problem}

% Notation shortcuts
\newcommand{\E}{\mathbb{E}}
\newcommand{\R}{\mathbb{R}}
\newcommand{\N}{\mathbb{N}}
\newcommand{\cS}{\mathcal{S}}
\newcommand{\cA}{\mathcal{A}}
\newcommand{\cR}{\mathcal{R}}
\newcommand{\cP}{\mathcal{P}}
\newcommand{\cZ}{\mathcal{Z}}
\newcommand{\cO}{\mathcal{O}}
\newcommand{\cI}{\mathcal{I}}
\newcommand{\cF}{\mathcal{F}}
\newcommand{\cM}{\mathcal{M}}
\newcommand{\cD}{\mathcal{D}}
\newcommand{\cT}{\mathcal{T}}
\newcommand{\cB}{\mathcal{B}}
\newcommand{\cH}{\mathcal{H}}
\newcommand{\Vpi}{V^{\pi}}
\newcommand{\Qpi}{Q^{\pi}}
\newcommand{\Vstar}{V^{*}}
\newcommand{\Qstar}{Q^{*}}
\newcommand{\pistar}{\pi^{*}}
\DeclareMathOperator*{\argmax}{arg\,max}
\DeclareMathOperator*{\argmin}{arg\,min}
\DeclareMathOperator{\KL}{KL}
\DeclareMathOperator{\supp}{supp}
\DeclareMathOperator{\Regret}{Regret}
\DeclareMathOperator{\Span}{span}

\title{\textbf{Formal Foundations of Reinforcement Learning and Planning} \\[4pt]
\Large Mathematical Foundations for Agentic AI \\[8pt]
\normalsize Research Notes --- ETH Zurich ML Lecture Series}
\author{ETH Z\"urich --- Machine Learning Lecture Series}
\date{Spring 2026}

\begin{document}
\maketitle

\begin{abstract}
These notes develop the mathematical foundations of reinforcement learning (RL) and planning
that underpin modern agentic AI systems. Starting from the reward hypothesis and the
Bellman optimality principle, we build up to distributional RL, successor features,
information-directed sampling, and the options framework for temporal abstraction.
Throughout, we emphasize rigorous proofs---the contraction properties of
Bellman operators, distributional convergence in the Wasserstein metric,
and the Generalised Policy Improvement theorem---while connecting these
classical results to contemporary developments in large-scale agentic AI,
RLHF, test-time compute, and the ``bitter lesson.''
\end{abstract}

\tableofcontents
\newpage

%% =======================================================================
\section{The Reward Hypothesis and ``Reward is Enough''}
\label{sec:reward}
%% =======================================================================

\subsection{Statement and Formalization}

The \textbf{reward hypothesis} (Sutton, 2004) asserts:

\begin{tcolorbox}[colback=blue!5!white, colframe=blue!60!black, title=The Reward Hypothesis]
\textit{That all of what we mean by goals and purposes can be well thought of as the maximization of the expected value of the cumulative sum of a received scalar signal (called reward).}
\end{tcolorbox}

Silver et al.\ (2021) strengthened this to the ``Reward is Enough'' thesis: intelligence
and its associated abilities---knowledge, learning, perception, social intelligence,
language, generalization---arise as subgoals of reward maximization in a
\emph{sufficiently rich environment}.

\begin{definition}[Sufficiently Rich Environment]
\label{def:rich-env}
An environment $\cM = (\cS, \cA, P, R, \gamma)$ is \emph{sufficiently rich} for a capability
$\mathfrak{C}$ if every policy $\pi$ that achieves near-optimal expected return
$\E_\pi[\sum_{t=0}^\infty \gamma^t R_t] \geq V^* - \varepsilon$ necessarily exhibits
capability $\mathfrak{C}$.  Formally, there exists a predicate $\Phi_{\mathfrak{C}}$
on policies such that
\[
  V^\pi \geq V^* - \varepsilon \implies \Phi_{\mathfrak{C}}(\pi) = \text{true}.
\]
\end{definition}

The key insight is that the environment does the ``heavy lifting'': in a complex
enough world, achieving high reward \emph{requires} developing sophisticated
internal representations, planning capabilities, and even social reasoning.

\subsection{Connection to AIXI}

Hutter's AIXI framework (2004) formalizes the reward-maximization hypothesis in
its most general form. The AIXI agent selects actions to maximize expected future
reward under a universal prior:
\[
  a_t^{\text{AIXI}} = \argmax_{a_t} \sum_{o_t r_t} \cdots \max_{a_m} \sum_{o_m r_m}
    [r_t + \cdots + r_m] \sum_{\nu \in \cM_U} 2^{-K(\nu)} \nu(o_{1:m} r_{1:m} \mid a_{1:m}),
\]
where $K(\nu)$ is the Kolmogorov complexity of environment $\nu$, and $\cM_U$
is the class of all computable environments. AIXI is Pareto-optimal but
incomputable---it establishes reward maximization as \emph{sufficient} in
principle, while leaving the question of \emph{efficient} reward maximization
to practical RL.

\subsection{Counterarguments: Multi-Objective RL}

The reward hypothesis faces substantive objections:

\begin{enumerate}[label=(\roman*)]
  \item \textbf{Multi-objective trade-offs.} Real agents often face Pareto-incomparable
    objectives (e.g., safety vs.\ performance). While scalarization
    $R = \sum_i w_i R_i$ can in principle recover any Pareto-optimal policy, the
    \emph{choice} of weights $w$ itself requires a meta-objective
    (Roijers et al., 2013).
  \item \textbf{Non-Markovian rewards.} Many natural objectives (e.g., fairness
    over time) are not expressible as functions of the current state-action pair.
    Reward machines (Icarte et al., 2018) and temporal logic specifications
    partially address this.
  \item \textbf{Specification gaming.} Optimizing a proxy reward may diverge
    from the designer's true intent---the \emph{reward misspecification} problem
    central to AI alignment.
  \item \textbf{Bounded rationality.} Real agents have computational constraints
    that prevent exact reward maximization; satisficing or resource-bounded
    rationality may be more descriptive (Simon, 1955).
\end{enumerate}

Despite these objections, the reward hypothesis remains the dominant
paradigm and provides the mathematical scaffolding for RL theory.


%% =======================================================================
\section{Bellman Optimality and Dynamic Programming}
\label{sec:bellman}
%% =======================================================================

\subsection{MDP Formalism}

\begin{definition}[Markov Decision Process]
A (discounted) MDP is a tuple $\cM = (\cS, \cA, P, R, \gamma)$ where:
\begin{itemize}[nosep]
  \item $\cS$ is a (finite) state space, $\cA$ is a (finite) action space.
  \item $P : \cS \times \cA \to \Delta(\cS)$ is the transition kernel:
    $P(s' \mid s, a) = \Pr(S_{t+1} = s' \mid S_t = s, A_t = a)$.
  \item $R : \cS \times \cA \to \R$ is the expected reward function (or more
    generally, $R : \cS \times \cA \to \Delta(\R)$).
  \item $\gamma \in [0, 1)$ is the discount factor.
\end{itemize}
\end{definition}

\begin{definition}[Value Functions]
For a stationary policy $\pi : \cS \to \Delta(\cA)$:
\begin{align}
  \Vpi(s) &= \E_\pi\!\left[\sum_{t=0}^\infty \gamma^t R(S_t, A_t) \;\middle|\; S_0 = s\right], \label{eq:vpi} \\
  \Qpi(s, a) &= \E_\pi\!\left[\sum_{t=0}^\infty \gamma^t R(S_t, A_t) \;\middle|\; S_0 = s, A_0 = a\right]. \label{eq:qpi}
\end{align}
\end{definition}

\subsection{Bellman Equations}

The value functions satisfy the \emph{Bellman consistency equations}:
\begin{align}
  \Vpi(s) &= \sum_{a} \pi(a \mid s) \left[ R(s,a) + \gamma \sum_{s'} P(s' \mid s, a)\, \Vpi(s') \right], \label{eq:bellman-v}\\
  \Qpi(s,a) &= R(s,a) + \gamma \sum_{s'} P(s' \mid s, a) \sum_{a'} \pi(a' \mid s')\, \Qpi(s', a'). \label{eq:bellman-q}
\end{align}

The \emph{Bellman optimality equations} characterize the optimal value:
\begin{align}
  \Vstar(s) &= \max_{a \in \cA} \left[ R(s,a) + \gamma \sum_{s'} P(s' \mid s, a)\, \Vstar(s') \right], \label{eq:bellman-opt-v}\\
  \Qstar(s,a) &= R(s,a) + \gamma \sum_{s'} P(s' \mid s, a)\, \max_{a'} \Qstar(s', a'). \label{eq:bellman-opt-q}
\end{align}

\subsection{The Bellman Operator and Contraction}

\begin{definition}[Bellman Optimality Operator]
Define $\cT : \R^{|\cS|} \to \R^{|\cS|}$ by
\[
  (\cT V)(s) = \max_{a \in \cA} \left[ R(s,a) + \gamma \sum_{s'} P(s' \mid s, a)\, V(s') \right].
\]
\end{definition}

\begin{theorem}[$\gamma$-Contraction in Supremum Norm]
\label{thm:contraction}
The Bellman optimality operator $\cT$ is a $\gamma$-contraction in the
$\ell_\infty$ norm: for all $V, U \in \R^{|\cS|}$,
\[
  \|\cT V - \cT U\|_\infty \leq \gamma \|V - U\|_\infty.
\]
\end{theorem}

\begin{proof}
For any state $s \in \cS$, let $a^* = \argmax_a [R(s,a) + \gamma \sum_{s'} P(s'|s,a) V(s')]$. Then:
\begin{align*}
  (\cT V)(s) - (\cT U)(s)
  &= \max_a \bigl[ R(s,a) + \gamma \textstyle\sum_{s'} P(s'|s,a) V(s') \bigr]
     - \max_a \bigl[ R(s,a) + \gamma \textstyle\sum_{s'} P(s'|s,a) U(s') \bigr] \\
  &\leq R(s,a^*) + \gamma \textstyle\sum_{s'} P(s'|s,a^*) V(s')
     - R(s,a^*) - \gamma \textstyle\sum_{s'} P(s'|s,a^*) U(s') \\
  &= \gamma \textstyle\sum_{s'} P(s'|s,a^*) \bigl[ V(s') - U(s') \bigr] \\
  &\leq \gamma \textstyle\sum_{s'} P(s'|s,a^*)\, \|V - U\|_\infty \\
  &= \gamma \|V - U\|_\infty.
\end{align*}
By symmetry (swapping $V$ and $U$), we also obtain $(\cT U)(s) - (\cT V)(s) \leq \gamma\|V - U\|_\infty$.
Hence $|(\cT V)(s) - (\cT U)(s)| \leq \gamma\|V - U\|_\infty$ for all $s$,
and taking the supremum over $s$ yields the result.
\end{proof}

\begin{corollary}[Existence and Uniqueness of $V^*$]
\label{cor:vstar-unique}
By the Banach fixed-point theorem, since $(\R^{|\cS|}, \|\cdot\|_\infty)$ is a
complete metric space and $\cT$ is a $\gamma$-contraction, there exists a unique
$\Vstar \in \R^{|\cS|}$ satisfying $\cT \Vstar = \Vstar$. Moreover, value iteration
$V_{k+1} = \cT V_k$ converges geometrically: $\|V_k - \Vstar\|_\infty \leq \gamma^k \|V_0 - \Vstar\|_\infty$.
\end{corollary}

\subsection{The Deadly Triad}

Sutton and Barto (2018) identified the \textbf{deadly triad}: the combination of
\begin{enumerate}[nosep]
  \item \textbf{Function approximation} (representing $V$ or $Q$ with a parameterized function class),
  \item \textbf{Bootstrapping} (updating estimates from other estimates, as in TD learning),
  \item \textbf{Off-policy learning} (learning about a policy different from the behavior policy),
\end{enumerate}
which can cause \emph{divergence} of the learning process. Any two of the three
are safe; the combination of all three can be unstable.

\begin{theorem}[Linear TD Convergence, Tsitsiklis and Van Roy, 1997]
\label{thm:linear-td}
Consider on-policy TD(0) with linear function approximation $\hat{V}(s; \mathbf{w}) = \boldsymbol{\phi}(s)^\top \mathbf{w}$,
where $\boldsymbol{\phi}(s) \in \R^d$ is a feature vector. Let $\mathbf{\Phi} \in \R^{|\cS| \times d}$ be the feature
matrix and $\mathbf{D} = \mathrm{diag}(d_\pi(s))$ the diagonal matrix of the stationary distribution.
If $\mathbf{\Phi}$ has full column rank, then the TD fixed point
\[
  \mathbf{w}^* = (\mathbf{\Phi}^\top \mathbf{D} (\mathbf{\Phi} - \gamma \mathbf{P}_\pi \mathbf{\Phi}))^{-1} \mathbf{\Phi}^\top \mathbf{D}\, \mathbf{r}_\pi
\]
exists and is unique, and TD(0) converges to $\mathbf{w}^*$ with probability 1 under
standard stochastic approximation conditions (Robbins--Monro step sizes). The
approximation error satisfies:
\[
  \|\mathbf{\Phi} \mathbf{w}^* - \Vpi\|_{\mathbf{D}} \leq \frac{1}{\sqrt{1 - \gamma^2}} \min_{\mathbf{w}} \|\mathbf{\Phi} \mathbf{w} - \Vpi\|_{\mathbf{D}}.
\]
\end{theorem}


%% =======================================================================
\section{Distributional Reinforcement Learning}
\label{sec:distributional}
%% =======================================================================

\subsection{Return Distributions}

Rather than modeling the \emph{expected} return, distributional RL (Bellemare, Dabney, and Munos, 2017)
models the full \emph{distribution} of the random return.

\begin{definition}[Return Distribution]
For a policy $\pi$, the random return from state $s$ taking action $a$ is
\[
  Z^\pi(s, a) = \sum_{t=0}^\infty \gamma^t R(S_t, A_t), \quad S_0 = s,\; A_0 = a,\; A_{t} \sim \pi(\cdot | S_t)\ \text{for } t \geq 1.
\]
The value function is recovered as $\Qpi(s,a) = \E[Z^\pi(s,a)]$.
\end{definition}

\subsection{Distributional Bellman Equation}

The return distribution satisfies a \emph{distributional} Bellman equation:
\[
  Z^\pi(s, a) \stackrel{D}{=} R(s, a) + \gamma\, Z^\pi(S', A'),
  \quad S' \sim P(\cdot \mid s, a),\; A' \sim \pi(\cdot \mid S').
\]
Here $\stackrel{D}{=}$ denotes equality in distribution. Define the
\textbf{distributional Bellman operator} $\cT^\pi$ acting on return-distribution
functions $Z : \cS \times \cA \to \cP(\R)$:
\[
  (\cT^\pi Z)(s,a) \stackrel{D}{=} R(s,a) + \gamma\, Z(S', A').
\]

\subsection{Contraction in Wasserstein Distance}

\begin{definition}[$p$-Wasserstein Distance]
For probability measures $\mu, \nu$ on $\R$ with finite $p$-th moments, the
$p$-Wasserstein distance is:
\[
  W_p(\mu, \nu) = \inf_{\text{coupling } (X,Y)} \bigl(\E[|X - Y|^p]\bigr)^{1/p},
\]
where the infimum is over all joint distributions of $(X, Y)$ with marginals $\mu$ and $\nu$.
For distributions on $\R$, via the quantile representation:
$W_p(\mu, \nu) = \bigl(\int_0^1 |F_\mu^{-1}(\tau) - F_\nu^{-1}(\tau)|^p \, d\tau\bigr)^{1/p}$.
\end{definition}

We lift the Wasserstein distance to return-distribution functions. For $Z_1, Z_2 : \cS \times \cA \to \cP(\R)$:
\[
  \bar{d}_p(Z_1, Z_2) = \sup_{s, a} W_p\bigl(Z_1(s,a),\, Z_2(s,a)\bigr).
\]

\begin{theorem}[Distributional Contraction, Bellemare et al., 2017]
\label{thm:dist-contraction}
The distributional Bellman operator $\cT^\pi$ is a $\gamma$-contraction in
$\bar{d}_p$ for any $p \geq 1$:
\[
  \bar{d}_p(\cT^\pi Z_1,\, \cT^\pi Z_2) \leq \gamma\, \bar{d}_p(Z_1, Z_2).
\]
\end{theorem}

\begin{proof}[Proof (coupling argument)]
Fix $(s, a) \in \cS \times \cA$. We have
\begin{align*}
  (\cT^\pi Z_1)(s,a) &\stackrel{D}{=} R(s,a) + \gamma\, Z_1(S', A'), \\
  (\cT^\pi Z_2)(s,a) &\stackrel{D}{=} R(s,a) + \gamma\, Z_2(S', A'),
\end{align*}
where $S' \sim P(\cdot|s,a)$ and $A' \sim \pi(\cdot|S')$.

We construct a coupling. For each $(s', a')$, let $(X_{s',a'}, Y_{s',a'})$ be an
\emph{optimal coupling} of $Z_1(s', a')$ and $Z_2(s', a')$, i.e.,
\[
  \E[|X_{s',a'} - Y_{s',a'}|^p] = W_p^p(Z_1(s',a'), Z_2(s',a')).
\]

Now couple $(\cT^\pi Z_1)(s,a)$ and $(\cT^\pi Z_2)(s,a)$ by using the \emph{same}
transition $(S', A')$ for both and using the optimal coupling for the downstream
returns. This yields random variables $\hat{X} = R(s,a) + \gamma X_{S',A'}$ and
$\hat{Y} = R(s,a) + \gamma Y_{S',A'}$ that are valid couplings.  Then:
\begin{align*}
  W_p^p\bigl((\cT^\pi Z_1)(s,a),\, (\cT^\pi Z_2)(s,a)\bigr)
  &\leq \E\bigl[|\hat{X} - \hat{Y}|^p\bigr] \\
  &= \E\bigl[|\gamma X_{S',A'} - \gamma Y_{S',A'}|^p\bigr] \\
  &= \gamma^p \,\E\bigl[|X_{S',A'} - Y_{S',A'}|^p\bigr] \\
  &= \gamma^p \sum_{s',a'} P(s'|s,a)\,\pi(a'|s')\, W_p^p(Z_1(s',a'), Z_2(s',a')) \\
  &\leq \gamma^p \,\bar{d}_p(Z_1, Z_2)^p.
\end{align*}
Taking $p$-th roots and the supremum over $(s, a)$:
\[
  \bar{d}_p(\cT^\pi Z_1, \cT^\pi Z_2) \leq \gamma\, \bar{d}_p(Z_1, Z_2). \qedhere
\]
\end{proof}

\begin{remark}
By the Banach fixed-point theorem applied to the complete metric space
$(\cZ, \bar{d}_p)$ (where $\cZ$ is the space of return-distribution functions with
finite $p$-th moments), the distributional Bellman equation has a unique
fixed point $Z^\pi$.  Crucially, the distributional Bellman \emph{optimality}
operator is \textbf{not} a contraction in any $\bar{d}_p$, which introduces
additional algorithmic challenges.
\end{remark}

\subsection{Algorithmic Instantiations}

\begin{itemize}[leftmargin=*]
  \item \textbf{C51} (Bellemare et al., 2017): Represents $Z(s,a)$ as a categorical
    distribution over a fixed set of $N = 51$ atoms $\{z_1, \ldots, z_N\}$.
    The distributional Bellman update is projected onto this support.

  \item \textbf{QR-DQN} (Dabney et al., 2018a): Parameterizes the return distribution
    by $N$ quantile values $\theta_1, \ldots, \theta_N$ at fixed quantile levels
    $\tau_i = (2i-1)/(2N)$. Minimizes the quantile Huber loss:
    \[
      \sum_{i=1}^N \sum_{j=1}^{N'} \rho_{\tau_i}^\kappa(\delta_{ij}),
      \quad \rho_\tau^\kappa(u) = |\tau - \mathbf{1}_{u < 0}| \cdot L_\kappa(u),
    \]
    where $L_\kappa$ is the Huber loss and $\delta_{ij} = r + \gamma \theta_j' - \theta_i$.

  \item \textbf{IQN} (Dabney et al., 2018b): Implicit Quantile Network---instead of
    fixed quantile levels, samples $\tau \sim \text{Uniform}(0,1)$ and learns the
    quantile function $F_Z^{-1}(\tau; s, a)$ directly using a neural network that
    takes $\tau$ as input.
\end{itemize}

\subsection{Neuroscience Connection}

\begin{tcolorbox}[colback=green!5!white, colframe=green!60!black, title=Distributional Coding in the Brain]
Dabney et al.\ (2020, \emph{Nature}) demonstrated that dopamine neurons in the
mouse brain encode not just reward prediction errors (as in classical TD
learning), but a \emph{distribution} of prediction errors. Different neurons
appear to encode different quantiles of the return distribution, forming a
population code for distributional value. This provides striking biological
evidence for the normative optimality of distributional RL: evolution arrived
at distributional representations independently.
\end{tcolorbox}


%% =======================================================================
\section{Successor Features and Generalised Policy Improvement}
\label{sec:sf}
%% =======================================================================

\subsection{Successor Representation}

The \textbf{successor representation} (Dayan, 1993) decomposes the value function
by separating \emph{what happens} (dynamics under $\pi$) from \emph{what is rewarding}
(the reward function).

Assume rewards are linear in state features: $R(s, a) = \boldsymbol{\phi}(s, a)^\top \mathbf{w}$,
where $\boldsymbol{\phi} : \cS \times \cA \to \R^d$ and $\mathbf{w} \in \R^d$.

\begin{definition}[Successor Features]
The successor features (SFs) under policy $\pi$ are:
\[
  \boldsymbol{\psi}^\pi(s, a) = \E_\pi\!\left[\sum_{t=0}^\infty \gamma^t \boldsymbol{\phi}(S_t, A_t) \;\middle|\; S_0 = s, A_0 = a\right] \in \R^d.
\]
\end{definition}

The $Q$-function decomposes as:
\begin{equation}
  \Qpi(s, a) = \boldsymbol{\psi}^\pi(s, a)^\top \mathbf{w}. \label{eq:sf-decomp}
\end{equation}

Critically, $\boldsymbol{\psi}^\pi$ depends only on $\pi$ and the dynamics $P$, not on $\mathbf{w}$.
Given a library of SFs $\{\boldsymbol{\psi}^{\pi_i}\}_{i=1}^n$ computed for different policies,
a new task specified by $\mathbf{w}'$ can be solved \emph{instantly}
by evaluating $\boldsymbol{\psi}^{\pi_i}(s,a)^\top \mathbf{w}'$ for each stored policy.

\subsection{Generalised Policy Improvement}

\begin{theorem}[Generalised Policy Improvement, Barreto et al., 2017]
\label{thm:gpi}
Let $\pi_1, \ldots, \pi_n$ be $n$ policies with known successor features
$\boldsymbol{\psi}^{\pi_1}, \ldots, \boldsymbol{\psi}^{\pi_n}$, and let $\mathbf{w}$ specify a new task with
$R(s,a) = \boldsymbol{\phi}(s,a)^\top \mathbf{w}$. Define the GPI policy:
\[
  \pi_{\mathrm{GPI}}(s) = \argmax_{a \in \cA} \max_{i \in \{1,\ldots,n\}} \boldsymbol{\psi}^{\pi_i}(s, a)^\top \mathbf{w}.
\]
Then $\pi_{\mathrm{GPI}}$ is at least as good as every $\pi_i$ on task $\mathbf{w}$:
\[
  Q^{\pi_{\mathrm{GPI}}}(s, a) \geq \max_{i} Q^{\pi_i}(s, a) \quad \forall\, s \in \cS,\; a \in \cA.
\]
In particular, $V^{\pi_{\mathrm{GPI}}}(s) \geq \max_i V^{\pi_i}(s)$ for all $s$.
\end{theorem}

\begin{proof}
We proceed by a policy improvement argument. For any $s \in \cS$:
\begin{align*}
  V^{\pi_{\mathrm{GPI}}}(s)
  &= \E_{\pi_{\mathrm{GPI}}}\!\left[\sum_{t=0}^\infty \gamma^t R(S_t, A_t) \;\middle|\; S_0 = s \right] \\
  &= \E_{\pi_{\mathrm{GPI}}}\!\left[R(s, A_0) + \gamma V^{\pi_{\mathrm{GPI}}}(S_1) \;\middle|\; S_0 = s\right].
\end{align*}
By construction, $\pi_{\mathrm{GPI}}$ selects $A_0 = \argmax_a \max_i \boldsymbol{\psi}^{\pi_i}(s,a)^\top \mathbf{w}$. Thus:
\[
  Q^{\pi_{\mathrm{GPI}}}(s, a) = R(s,a) + \gamma \E[V^{\pi_{\mathrm{GPI}}}(S_1) \mid S_0 = s, A_0 = a].
\]
We prove $Q^{\pi_{\mathrm{GPI}}}(s,a) \geq \max_i Q^{\pi_i}(s,a)$ by induction on the horizon.
For the base case ($\gamma = 0$), it is immediate since $Q^{\pi_i}(s,a) = R(s,a)$ for all $i$.

For the inductive step, assume $V^{\pi_{\mathrm{GPI}}}(s') \geq \max_i V^{\pi_i}(s')$ for
all $s'$ reachable in $k$ steps. Then for state $s$ at step $k+1$:
\begin{align*}
  Q^{\pi_{\mathrm{GPI}}}(s, a)
  &= R(s,a) + \gamma \sum_{s'} P(s'|s,a)\, V^{\pi_{\mathrm{GPI}}}(s') \\
  &\geq R(s,a) + \gamma \sum_{s'} P(s'|s,a) \max_i V^{\pi_i}(s') \\
  &\geq R(s,a) + \gamma \sum_{s'} P(s'|s,a)\, V^{\pi_i}(s') \quad \text{for all } i \\
  &= Q^{\pi_i}(s,a) \quad \text{for all } i.
\end{align*}
Taking the max over $i$ gives the result. The argument extends to the infinite-horizon
discounted case via the contraction mapping argument: let $Q_{\mathrm{GPI}}$ be the fixed
point of the operator $(\cT_{\mathrm{GPI}} Q)(s,a) = R(s,a) + \gamma \sum_{s'} P(s'|s,a) \max_a \max_i Q_i(s',a')$,
where we can show this operator dominates each $\cT^{\pi_i}$ pointwise.
\end{proof}

\subsection{Zero-Shot Transfer}

The GPI theorem enables \textbf{zero-shot transfer}: given successor features from
training tasks $\mathbf{w}_1, \ldots, \mathbf{w}_n$, performance on a new task $\mathbf{w}$ is bounded by:
\[
  V^*(s; \mathbf{w}) - V^{\pi_{\mathrm{GPI}}}(s; \mathbf{w}) \leq \frac{2\gamma}{1 - \gamma}
  \min_{i} \|\boldsymbol{\psi}^{\pi_i} - \boldsymbol{\psi}^*(s; \mathbf{w})\|_\infty \|\mathbf{w}\|_1.
\]
If the training tasks span the reward space well, this gap can be small even
for novel tasks.


%% =======================================================================
\section{Information-Directed Sampling}
\label{sec:ids}
%% =======================================================================

\subsection{The Exploration--Exploitation Dilemma}

The fundamental challenge in sequential decision-making is balancing
\emph{exploitation} (acting on current knowledge) with \emph{exploration}
(gathering information to improve future decisions). Classical approaches include:

\begin{itemize}[nosep]
  \item \textbf{UCB (Upper Confidence Bound)}: Act optimistically---select actions whose
    confidence-upper-bound on value is highest. Achieves $O(\sqrt{KT \log T})$
    regret for $K$-armed bandits (Auer et al., 2002).
  \item \textbf{Thompson Sampling}: Sample a model from the posterior and act
    greedily with respect to the sample. Achieves $O(\sqrt{KT \log K})$
    Bayesian regret (Russo and Van Roy, 2014).
\end{itemize}

Both can be suboptimal in structured problems. \textbf{Information-Directed Sampling}
(IDS), due to Russo and Van Roy (2014, 2018), provides a principled alternative.

\subsection{The Information Ratio}

\begin{definition}[Information Ratio]
Let $\pi_t$ be a (possibly randomized) action-selection rule at time $t$.
Given posterior $P_t$ over the environment, define:
\begin{itemize}[nosep]
  \item \textbf{Expected instantaneous regret}: $\Delta_t(\pi_t) = \E_{A_t \sim \pi_t}[\E_{P_t}[R(A^*) - R(A_t)]]$,
    where $A^*$ is the optimal action under the true model.
  \item \textbf{Information gain}: $g_t(\pi_t) = I_{P_t}(\theta;\, (A_t, O_t) \mid A_t \sim \pi_t)$,
    the mutual information between the unknown parameter $\theta$ and the
    observation $O_t$ under action $A_t$.
\end{itemize}
The \textbf{information ratio} is:
\begin{equation}
  \Gamma_t(\pi_t) = \frac{\Delta_t(\pi_t)^2}{g_t(\pi_t)}. \label{eq:info-ratio}
\end{equation}
\end{definition}

\begin{definition}[Information-Directed Sampling]
The IDS policy minimizes the information ratio:
\[
  \pi_t^{\mathrm{IDS}} = \argmin_{\pi_t \in \Delta(\cA)} \Gamma_t(\pi_t)
  = \argmin_{\pi_t} \frac{\Delta_t(\pi_t)^2}{g_t(\pi_t)}.
\]
\end{definition}

This trades off regret and information optimally: IDS takes actions that yield the
best ratio of ``regret squared per bit of information gained.''

\subsection{Regret Bound}

\begin{theorem}[IDS Regret Bound, Russo and Van Roy, 2014]
\label{thm:ids-regret}
If $\Gamma_t(\pi_t^{\mathrm{IDS}}) \leq \bar{\Gamma}$ for all $t$ (uniformly bounded
information ratio), and $H(\theta) < \infty$ is the entropy of the prior, then the
Bayesian regret of IDS satisfies:
\[
  \Regret(T) = \E\!\left[\sum_{t=1}^T \Delta_t\right] \leq \sqrt{\bar{\Gamma} \cdot H(\theta) \cdot T}.
\]
\end{theorem}

\begin{proof}[Proof sketch]
By the definition of the information ratio:
$\Delta_t^2 \leq \bar{\Gamma} \cdot g_t$ for all $t$.
Summing over $t$ and applying Jensen's inequality:
\[
  \left(\frac{1}{T}\sum_{t=1}^T \Delta_t\right)^2 \leq \frac{1}{T}\sum_{t=1}^T \Delta_t^2
  \leq \frac{\bar{\Gamma}}{T} \sum_{t=1}^T g_t.
\]
The key insight is the \emph{information-theoretic chain rule}: the total information
gained about $\theta$ cannot exceed the prior entropy,
$\sum_{t=1}^T g_t \leq H(\theta)$.
This is because $I(\theta; O_{1:T}) \leq H(\theta)$, and the chain rule for
mutual information gives $I(\theta; O_{1:T}) = \sum_t I(\theta; O_t \mid O_{1:t-1}) = \sum_t g_t$
(adapting for the adaptive action selection). Hence:
\[
  \Regret(T) = \sum_{t=1}^T \Delta_t \leq \sqrt{T \sum_{t=1}^T \Delta_t^2}
  \leq \sqrt{T \cdot \bar{\Gamma} \cdot H(\theta)}. \qedhere
\]
\end{proof}

\begin{remark}
For $K$-armed Bernoulli bandits, IDS achieves $\bar{\Gamma} \leq K/2$, yielding
$O(\sqrt{KT \log K})$ Bayesian regret (matching Thompson sampling). For linear
bandits in $d$ dimensions, $\bar{\Gamma} \leq d/2$, giving $O(d\sqrt{T})$ regret
that can \emph{improve} upon both UCB and Thompson sampling in structured problems.
\end{remark}

\subsection{Comparison and Connection to Active Inference}

\begin{center}
\begin{tabular}{lccc}
  \toprule
  \textbf{Method} & \textbf{Principle} & \textbf{Regret ($K$-armed)} & \textbf{Structured?} \\
  \midrule
  UCB & Optimism & $O(\sqrt{KT \log T})$ & Moderate \\
  Thompson Sampling & Posterior sampling & $O(\sqrt{KT \log K})$ & Good \\
  IDS & Min info ratio & $O(\sqrt{KT \log K})$ & Excellent \\
  \bottomrule
\end{tabular}
\end{center}

IDS is closely related to the \textbf{active inference} framework from neuroscience
(Friston, 2010), which postulates that biological agents minimize \emph{expected free energy}---a
quantity that decomposes into an exploitation term (expected reward or pragmatic value)
and an exploration term (expected information gain or epistemic value). IDS can be
viewed as a decision-theoretic formalization of this decomposition.


%% =======================================================================
\section{The Options Framework}
\label{sec:options}
%% =======================================================================

\subsection{Temporal Abstraction}

Primitive actions operate at the finest time scale. Many tasks are better
described at multiple temporal scales: ``navigate to the kitchen'' is a
temporally extended action composed of many primitive steps. The
\textbf{options framework} (Sutton, Precup, and Singh, 1999) formalizes
this via \emph{temporal abstraction}.

\begin{definition}[Option]
An option $o = (\cI_o, \pi_o, \beta_o)$ consists of:
\begin{itemize}[nosep]
  \item $\cI_o \subseteq \cS$: the \textbf{initiation set} (states where $o$ can begin),
  \item $\pi_o : \cS \to \Delta(\cA)$: the \textbf{intra-option policy},
  \item $\beta_o : \cS \to [0, 1]$: the \textbf{termination function}
    ($\beta_o(s)$ is the probability of terminating in state $s$).
\end{itemize}
\end{definition}

Primitive actions are options with $\cI_o = \cS$, $\pi_o(\cdot|s) = \mathbf{1}_a$,
and $\beta_o \equiv 1$ (immediate termination).

\subsection{Semi-MDP Formulation}

When executing options, the time between decisions is variable. This gives
rise to a \textbf{semi-MDP} (SMDP) over options. For a policy-over-options
$\mu : \cS \to \Delta(\cO)$ (where $\cO$ is the set of available options),
the \textbf{option-value function} satisfies the SMDP Bellman equation:
\begin{equation}
  Q_\cO^\mu(s, o) = \E\!\left[\sum_{k=0}^{\tau_o - 1} \gamma^k R_{t+k} + \gamma^{\tau_o}
  \sum_{o'} \mu(o' | S_{t+\tau_o})\, Q_\cO^\mu(S_{t+\tau_o}, o') \;\middle|\; S_t = s, O_t = o \right],
  \label{eq:option-bellman}
\end{equation}
where $\tau_o$ is the (random) duration of option $o$.

With the \emph{call-and-return} option execution model, this can be written
more compactly using the \textbf{option-conditional transition model}:
\[
  P_o(s', k \mid s) = \Pr(S_{t+k} = s', \tau_o = k \mid S_t = s),
\]
and the expected discounted reward during execution:
\[
  r_o(s) = \E\!\left[\sum_{k=0}^{\tau_o - 1} \gamma^k R_{t+k} \;\middle|\; S_t = s \right].
\]

\subsection{Eigenoptions and Option Discovery}

A key challenge is \emph{discovering} useful options automatically.
\textbf{Eigenoptions} (Machado et al., 2017) uses the graph Laplacian of the
state-transition graph: eigenvectors of the Laplacian identify bottleneck
states and natural sub-goals.

Given the graph Laplacian $\mathbf{L} = \mathbf{D}_G - \mathbf{A}$ (where $\mathbf{A}$ is the adjacency
matrix and $\mathbf{D}_G$ the degree matrix), the Fiedler vector (second-smallest
eigenvector of $\mathbf{L}$) partitions the state space into two weakly connected
components. Options that navigate between these components correspond to
traversing bottleneck states and are useful for exploration.

More generally, the $k$-th eigenvector $\mathbf{e}_k$ of $\mathbf{L}$ defines an intrinsic
reward $r_k(s, s') = \mathbf{e}_k(s') - \mathbf{e}_k(s)$, and the corresponding eigenoption
is the policy maximizing cumulative $r_k$.

\subsection{Option-Critic Architecture}

The \textbf{Option-Critic} (Bacon, Harb, and Precup, 2017) learns options
end-to-end using the policy gradient theorem extended to options:

\begin{align}
  \frac{\partial V_\cO(s)}{\partial \theta_{\pi_o}} &\propto \sum_{s,o} \mu_\cO(s, o \mid s_0)
  \sum_a \frac{\partial \pi_{o}(a|s)}{\partial \theta_{\pi_o}} Q_U(s, o, a), \\
  \frac{\partial V_\cO(s)}{\partial \theta_{\beta_o}} &\propto -\sum_{s',o} \mu_\cO(s', o \mid s_0)
  \frac{\partial \beta_o(s')}{\partial \theta_{\beta_o}} A_\cO(s', o),
\end{align}
where $Q_U(s, o, a)$ is the value of executing action $a$ within option $o$ at state $s$,
and $A_\cO(s', o) = Q_\cO(s', o) - V_\cO(s')$ is the advantage of continuing option $o$.

\subsection{Tool Use as Options}

In agentic AI systems, \textbf{tool use} maps naturally to the options framework:
\begin{itemize}[nosep]
  \item A \emph{tool} (web search, code execution, calculator) is an option with
    its own internal policy (the tool's execution logic).
  \item The \emph{initiation set} corresponds to contexts where the tool is applicable.
  \item The \emph{termination condition} is when the tool returns a result.
  \item The agent's \emph{policy over options} is the decision of which tool to invoke.
\end{itemize}
This abstraction allows treating tool-augmented agents within the standard
RL/planning framework (see Section~\ref{sec:agentic}).


%% =======================================================================
\section{The Bitter Lesson and Scaling Laws}
\label{sec:bitter}
%% =======================================================================

\subsection{The Bitter Lesson}

Sutton's ``Bitter Lesson'' (2019) observes that, over the history of AI research,
methods that \emph{leverage computation}---particularly search and learning---have
consistently outperformed methods that \emph{encode human knowledge}. This lesson
is ``bitter'' because researchers repeatedly invest effort in domain-specific
engineering, only to be surpassed by general methods that scale with compute.

The key examples span decades: chess (Deep Blue to AlphaZero), Go (handcrafted
features to AlphaGo/AlphaZero), speech recognition (HMMs to deep learning),
computer vision (SIFT/HOG to CNNs), and NLP (parse trees to Transformers).

\subsection{Scaling Laws for RL}

Scaling laws in supervised learning (Kaplan et al., 2020; Hoffmann et al., 2022)
describe power-law relationships between loss and compute/data/parameters.
Analogous scaling phenomena appear in RL:

\begin{enumerate}[nosep]
  \item \textbf{Sample complexity}: For tabular MDPs, PAC-MDP bounds scale as
    $\tilde{O}(|\cS|^2|\cA|/\varepsilon^2)$ (Strehl et al., 2009). With function
    approximation, this depends on the effective dimension of the function class.
  \item \textbf{Compute scaling in games}: AlphaZero's Elo rating scales
    log-linearly with training compute. MuZero exhibits similar scaling.
  \item \textbf{Test-time compute scaling}: In reasoning tasks, increasing
    search budget at inference time (e.g., more MCTS rollouts) yields predictable
    performance improvements following power-law curves.
\end{enumerate}

\subsection{The Era of Experience}

Silver, Sutton, and collaborators (2025) argue for an ``Era of Experience''
in which AI agents learn primarily from \emph{self-generated experience}
rather than human-curated data. Key claims include:

\begin{itemize}[nosep]
  \item The data bottleneck for LLMs (limited human text) can be overcome by
    agents generating their own training signal through interaction.
  \item The combination of world models, planning, and self-play enables
    superhuman performance without human demonstrations.
  \item This extends the bitter lesson: not only computation, but \emph{experience}
    should be scaled, and the most general-purpose approach is to let agents
    learn from their own interaction data.
\end{itemize}


%% =======================================================================
\section{Connections to Modern Agentic AI}
\label{sec:agentic}
%% =======================================================================

The RL foundations developed above connect directly to the architecture and
training of modern agentic AI systems.

\subsection{RLHF and RLAIF}

\textbf{Reinforcement Learning from Human Feedback} (RLHF; Christiano et al., 2017;
Ouyang et al., 2022) trains a reward model from human preference comparisons,
then optimizes a language model policy against this reward using RL (typically PPO).
The objective is:
\[
  \max_\pi \E_{x \sim \cD,\, y \sim \pi(\cdot|x)}\bigl[r_\phi(x, y)\bigr]
  - \beta\, \KL\!\bigl[\pi(\cdot|x) \,\|\, \pi_{\mathrm{ref}}(\cdot|x)\bigr],
\]
where $r_\phi$ is the learned reward model and the KL penalty prevents
collapse away from the reference policy $\pi_{\mathrm{ref}}$.

\textbf{RLAIF} (Bai et al., 2022) replaces human preferences with AI-generated
preferences, connecting to the ``Era of Experience'' thesis: the agent generates
its own training signal.

The closed-form solution of the KL-constrained problem is:
\[
  \pi^*(y|x) = \frac{1}{Z(x)} \pi_{\mathrm{ref}}(y|x) \exp\!\left(\frac{r_\phi(x,y)}{\beta}\right),
\]
which motivates Direct Preference Optimization (DPO; Rafailov et al., 2023).

\subsection{In-Context Learning as Implicit RL}

Transformer-based models exhibit \emph{in-context learning}: they adapt their
behavior based on examples in the prompt without weight updates. This can be
viewed as \emph{implicit} RL or Bayesian inference:

\begin{itemize}[nosep]
  \item The prompt context is the ``trajectory'' so far.
  \item The model's next-token prediction implements a policy conditioned on history.
  \item Transformers trained on diverse tasks can implement in-context
    algorithms that approximate Bayes-optimal policies (M\"uller et al., 2022;
    Garg et al., 2022).
  \item Algorithm distillation (Laskin et al., 2022) explicitly trains
    Transformers on RL learning histories, enabling in-context RL.
\end{itemize}

\subsection{Test-Time Compute as MCTS}

Recent work on ``thinking'' models (e.g., chain-of-thought, tree-of-thought)
allocates additional computation at inference time to solve harder problems.
This is directly analogous to MCTS:

\begin{itemize}[nosep]
  \item \textbf{Chain-of-thought}: A single rollout (analogous to a single MCTS simulation).
  \item \textbf{Tree-of-thought}: Multiple branching rollouts with backtracking
    (explicit tree search).
  \item \textbf{Best-of-$N$ sampling}: Sample $N$ complete solutions and select
    the best (analogous to root parallelization in MCTS).
  \item The value function in MCTS corresponds to the model's self-evaluation
    capability (process reward models).
\end{itemize}

\subsection{Multi-Step Tasks as POMDPs}

Agentic tasks (web browsing, coding, research) are naturally modeled as
\textbf{POMDPs}: the agent observes partial information (screen content, file
contents, API responses) and must maintain beliefs about the hidden state
(codebase structure, user intent, task progress). The agent's ``scratchpad''
or ``memory'' serves as a belief state, and tool use provides observations
that reduce uncertainty.

\subsection{Multi-Agent RL and LLM Agents}

Littman's framework of \textbf{Markov games} (1994) extends MDPs to multiple agents:
\begin{definition}[Markov Game]
A Markov game for $n$ agents is $(\cS, \cA_1, \ldots, \cA_n, P, R_1, \ldots, R_n, \gamma)$,
where $P : \cS \times \cA_1 \times \cdots \times \cA_n \to \Delta(\cS)$ and each
$R_i$ is agent $i$'s reward function.
\end{definition}

Multi-agent LLM systems (e.g., debate, negotiation, collaborative coding) are
instances of Markov games. The solution concept shifts from optimality to
\emph{equilibria} (Nash, correlated, or coarse-correlated equilibria), and
convergence guarantees become substantially more challenging.


%% =======================================================================
\section{Open Problems}
\label{sec:open}
%% =======================================================================

\begin{openproblem}[Non-Stationarity in RLHF]
Human preferences evolve over time, and reward models trained on past preferences
may become stale. How should RL agents adapt to non-stationary reward functions
while maintaining alignment guarantees? Connections to restless bandits
(Whittle, 1988) and non-stationary MDPs remain underexplored.
\end{openproblem}

\begin{openproblem}[Sample Complexity of Distributional RL]
While distributional RL empirically outperforms expectation-based methods,
the theoretical sample complexity advantages are not well understood.
Under what conditions does learning the full return distribution provably
improve policy optimization, beyond the benefits of auxiliary representation learning?
\end{openproblem}

\begin{openproblem}[Scalable Option Discovery]
Current option discovery methods (eigenoptions, interest functions, information-theoretic
objectives) struggle to scale to high-dimensional, continuous state spaces
encountered in real-world agentic tasks. Can foundation models serve as
priors for option discovery, identifying meaningful temporal abstractions
from language descriptions?
\end{openproblem}

\begin{openproblem}[Theoretical Foundations of Test-Time Compute Scaling]
Empirical scaling laws for test-time compute (more ``thinking'' = better performance)
lack rigorous theoretical grounding. What are the computational complexity-theoretic
characterizations of problems where test-time search provably helps?
How does the value of additional computation decay, and what is the optimal
compute allocation between training and inference?
\end{openproblem}

\begin{openproblem}[Credit Assignment in Long-Horizon Agentic Tasks]
Agentic tasks (multi-step coding, research, web navigation) involve hundreds or
thousands of steps with sparse, delayed rewards. The credit assignment problem---which
actions were responsible for success or failure---is exacerbated at this scale.
Existing methods (eligibility traces, hindsight relabeling, attention-based credit)
do not scale gracefully. Can successor features or distributional methods provide
better credit assignment in these regimes?
\end{openproblem}

\begin{openproblem}[Safe Exploration in Open-Ended Environments]
Real-world agentic AI systems must explore safely---avoiding irreversible
actions (deleting files, sending emails, financial transactions) while still
learning effectively. The constrained MDP framework provides some tools, but
the constraint set itself is often unknown and must be learned. How can
IDS-style information-theoretic exploration be combined with safety constraints?
\end{openproblem}

\begin{openproblem}[Bridging Language and Formal Reward Specification]
Natural language instructions are the primary interface for agentic AI, yet
translating them into formal reward functions remains error-prone. Reward
machines (Icarte et al., 2018) and linear temporal logic offer structured
approaches, but grounding language to these formalisms at scale is unsolved.
Can LLMs themselves serve as reward specification languages, and if so,
what are the alignment implications?
\end{openproblem}


%% =======================================================================
\section{Annotated Bibliography}
\label{sec:bib}
%% =======================================================================

\begin{enumerate}[label={[\arabic*]}, leftmargin=*, itemsep=6pt]

\item \textbf{Sutton, R.~S. and Barto, A.~G.} (2018).
  \textit{Reinforcement Learning: An Introduction} (2nd ed.). MIT Press. \\
  The definitive textbook. Chapters 3--4 cover MDPs and dynamic programming;
  Chapters 9--11 cover function approximation including the deadly triad;
  Chapter 12 covers eligibility traces.

\item \textbf{Silver, D., Singh, S., Precup, D., and Sutton, R.~S.} (2021).
  ``Reward is Enough.'' \textit{Artificial Intelligence}, 299, 103535. \\
  Argues that reward maximization in sufficiently complex environments is
  sufficient to give rise to all facets of intelligence. Provocative and
  foundational for the reward hypothesis.

\item \textbf{Bellemare, M.~G., Dabney, W., and Munos, R.} (2017).
  ``A Distributional Perspective on Reinforcement Learning.''
  \textit{Proceedings of ICML}. \\
  Introduces distributional RL, proves the distributional Bellman operator
  is a contraction in the Wasserstein metric, and proposes C51.

\item \textbf{Dabney, W., Ostrovski, G., Silver, D., and Munos, R.} (2018a).
  ``Implicit Quantile Networks for Distributional Reinforcement Learning.''
  \textit{Proceedings of ICML}. \\
  Introduces IQN, which learns the full quantile function using implicit
  representations. Achieves state-of-the-art on Atari.

\item \textbf{Dabney, W., Kurth-Nelson, Z., et al.} (2020).
  ``A distributional code for value in dopamine-based reinforcement learning.''
  \textit{Nature}, 577, 671--675. \\
  Landmark neuroscience paper showing dopamine neurons encode distributional
  reward prediction errors, vindicating distributional RL as normatively
  and descriptively accurate.

\item \textbf{Barreto, A., Dabney, W., Munos, R., et al.} (2017).
  ``Successor Features for Transfer in Reinforcement Learning.''
  \textit{Proceedings of NeurIPS}. \\
  Formalizes successor features and proves the Generalised Policy Improvement
  theorem. Demonstrates zero-shot transfer across tasks.

\item \textbf{Dayan, P.} (1993).
  ``Improving Generalization for Temporal Difference Learning: The Successor
  Representation.'' \textit{Neural Computation}, 5(4), 613--624. \\
  The original successor representation paper. Separates environment dynamics
  from reward, enabling efficient recomputation when rewards change.

\item \textbf{Russo, D. and Van Roy, B.} (2014).
  ``Learning to Optimize via Information-Directed Sampling.''
  \textit{Proceedings of NeurIPS}. \\
  Introduces IDS and the information ratio. Proves Bayesian regret bounds
  that can improve upon UCB and Thompson sampling.

\item \textbf{Russo, D. and Van Roy, B.} (2018).
  ``Learning to Optimize via Information-Directed Sampling.''
  \textit{Operations Research}, 66(1), 230--252. \\
  Journal version with extended analysis, including linear bandits and
  connections to rate-distortion theory.

\item \textbf{Sutton, R.~S., Precup, D., and Singh, S.} (1999).
  ``Between MDPs and semi-MDPs: A Framework for Temporal Abstraction in
  Reinforcement Learning.'' \textit{Artificial Intelligence}, 112, 181--211. \\
  The foundational options paper. Defines options, proves convergence of
  option-level value iteration, and establishes the SMDP connection.

\item \textbf{Bacon, P.-L., Harb, J., and Precup, D.} (2017).
  ``The Option-Critic Architecture.''
  \textit{Proceedings of AAAI}. \\
  Derives policy gradient theorems for learning options end-to-end,
  including both intra-option policies and termination conditions.

\item \textbf{Szepesv\'ari, C.} (2010).
  \textit{Algorithms for Reinforcement Learning}. Morgan \& Claypool. \\
  Concise monograph with rigorous treatment of TD learning, policy gradient
  methods, and actor-critic algorithms. Excellent for mathematical foundations.

\item \textbf{Littman, M.~L.} (1994).
  ``Markov Games as a Framework for Multi-Agent Reinforcement Learning.''
  \textit{Proceedings of ICML}. \\
  Extends MDPs to multi-agent settings. Introduces minimax-Q for zero-sum
  games and friend-or-foe Q-learning.

\item \textbf{Hutter, M.} (2004).
  \textit{Universal Artificial Intelligence: Sequential Decisions Based on
  Algorithmic Probability}. Springer. \\
  Defines AIXI, the Bayes-optimal agent over all computable environments.
  Establishes theoretical limits of reward maximization.

\item \textbf{Sutton, R.~S.} (2019).
  ``The Bitter Lesson.'' Blog post. \\
  Influential essay arguing that compute-leveraging methods (search and
  learning) consistently outperform knowledge-engineering approaches.

\item \textbf{Silver, D., Sutton, R.~S., et al.} (2025).
  ``Welcome to the Era of Experience.'' Preprint. \\
  Argues that AI should transition from learning on human-generated data to
  learning from self-generated experience via RL and world models.

\item \textbf{Ouyang, L., Wu, J., Jiang, X., et al.} (2022).
  ``Training language models to follow instructions with human feedback.''
  \textit{Proceedings of NeurIPS}. \\
  The InstructGPT paper. Demonstrates RLHF at scale: reward modeling from
  human comparisons followed by PPO optimization. Foundation for ChatGPT.

\item \textbf{Rafailov, R., Sharma, A., Mitchell, E., et al.} (2023).
  ``Direct Preference Optimization: Your Language Model is Secretly a Reward
  Model.'' \textit{Proceedings of NeurIPS}. \\
  Derives DPO from the closed-form solution of the KL-constrained RLHF
  objective, eliminating the need for explicit reward modeling.

\item \textbf{Munos, R., Szepesv\'ari, C., and Audibert, J.-Y.} (2008).
  ``Finite-Time Bounds for Fitted Value Iteration.''
  \textit{Journal of Machine Learning Research}, 9, 815--857. \\
  Rigorous finite-sample analysis of approximate dynamic programming with
  function approximation. Establishes error propagation bounds that depend
  on the inherent Bellman error and concentrability coefficients.

\item \textbf{Machado, M.~C., Rosenbaum, C., Guo, X., et al.} (2017).
  ``Eigenoption Discovery through the Deep Successor Representation.''
  \textit{Proceedings of ICLR}. \\
  Connects the graph Laplacian of the state space to option discovery via
  eigenvectors, enabling principled construction of temporally extended actions.

\end{enumerate}


%% =======================================================================
\section*{Acknowledgments}

These notes draw on the cumulative contributions of the RL research community.
We particularly acknowledge the foundational work of Richard Sutton, Andrew Barto,
Csaba Szepesv\'ari, and the DeepMind RL group. Any errors are our own.

\end{document}
